% چکیده فارسی
\chapter*{چکیده}
\phantomsection
\addcontentsline{toc}{chapter}{چکیده}

حل عددی معادلات دیفرانسیل جزئی غیرخطی از مسائل بنیادی محاسبات علمی است.
روش‌های عددی کلاسیک هرچند دقیق و پرکاربردند، به شبکه محاسباتی وابسته‌اند و
در ابعاد بالا با مشکل مواجه می‌شوند؛ از سوی دیگر، مدل‌های یادگیری عمیقِ
صرفاً داده‌محور به حجم زیادی داده آموزشی نیاز دارند و سازگاری خروجی آن‌ها
با قوانین فیزیک تضمین‌شده نیست. در سال‌های اخیر، چارچوب شبکه‌های عصبی آگاه
از فیزیک به‌عنوان راهی برای ترکیب دانش فیزیکی (معادله حاکم بر سیستم) با
یادگیری ماشین معرفی شده است. هدف این پایان‌نامه، مطالعه، پیاده‌سازی و
ارزیابی این چارچوب برای حل مسائل مستقیم و معکوس معادلات دیفرانسیل جزئی
غیرخطی است.

در این پژوهش، ابتدا مبانی نظری لازم شامل معادلات دیفرانسیل جزئی، روش‌های
عددی کلاسیک، شبکه‌های عصبی، مشتق‌گیری خودکار و پیشینه یادگیری ماشین آگاه
از فیزیک مرور شد. سپس چارچوب شبکه آگاه از فیزیک شامل تعریف شبکه باقیمانده،
تابع هزینه ترکیبی داده و فیزیک، و مدل‌های زمان‌پیوسته و زمان‌گسسته تشریح
گردید. در بخش عملی، مدل زمان‌پیوسته برای معادله برگرز یک‌بعدی با زبان
پایتون و کتابخانه جکس پیاده‌سازی شد و یک حل‌گر مرجع مستقل بر پایه روش طیفی
فوریه و گام‌زنی زمانی مرتبه چهار برای تولید جواب دقیق توسعه یافت و با جواب
نیمه‌تحلیلی کول ـ هوپف راستی‌آزمایی شد.

در مسئله مستقیم، جواب کل دامنه زمانی ـ مکانی تنها با ۲۰۰ نقطه داده اولیه
و مرزی بازسازی شد و خطای نسبی \FwdRelErr\ در نرم $L_2$ به دست آمد؛ خطا
عمدتاً در لایه باریک جبهه موج متمرکز بود. در مسئله معکوس، ضرایب معادله از
روی ۵۰۰۰ مشاهده پراکنده شناسایی شدند: ضریب انتقال با خطای \LamIErrClean\ و
ضریب پخش با خطای \LamIIErrClean\ در حالت بدون نویز، و به‌ترتیب با خطای
\LamIErrNoisy\ و \LamIIErrNoisy\ در حضور یک درصد نویز. نتایج نشان می‌دهد که
این چارچوب با داده اندک به جوابی قابل‌قبول می‌رسد، در برابر نویز پایداری
قابل‌قبولی دارد و گلوگاه اصلی آن، تفکیک لایه‌های تند جواب و حساسیت ذاتی
شناسایی ضرایب کوچک است.

\textbf{کلیدواژه‌ها:} شبکه عصبی آگاه از فیزیک، معادله دیفرانسیل جزئی
غیرخطی، مسئله مستقیم، مسئله معکوس، معادله برگرز، مشتق‌گیری خودکار
